\documentclass[10pt,a4paper]{amsart}
\usepackage[margin=19mm]{geometry}
\usepackage{amsmath,amssymb,mathtools}
\usepackage[scr=rsfs,cal=euler]{mathalfa}
\usepackage{xfrac}
\usepackage[unicode,hidelinks,bookmarks=false]{hyperref}
\newcommand{\ie}{i.e.\ }
\newcommand{\cf}{cf.\ }
\newcommand{\resp}{respectively\ }
\newcommand{\Subsection}{\subsection}
\providecommand{\nobreakdash}{\nobreak\hbox{-}}
\setlength{\emergencystretch}{2em}
\multlinegap0pt
\usepackage{bm}
\usepackage{bbm}
\usepackage{dsfont}
\usepackage{xspace}
\usepackage{cancel}
\usepackage{mathtools}
\usepackage{amsthm}
\makeatletter
\@ifundefined{lem}{\newtheorem{lem}{Lemma}}{}
\makeatother
\title[Combinatorial Reduction of Set Functions and Matroid Permuta]{Combinatorial Reduction of Set Functions and Matroid Permutations through Minor Product Assignment}
\author{Mario Angelelli}
\date{Source: arXiv:2107.01038v15 (CC BY 4.0)}
\begin{document}
\maketitle

\noindent\emph{Source and licence.} Combinatorial Reduction of Set Functions and Matroid Permutations through Minor Product Assignment. Version arXiv:2107.01038v15 (CC BY 4.0), distributed under the Creative Commons Attribution 4.0 International licence, as recorded in the arXiv licence metadata for this exact version and shown on the abstract page https://arxiv.org/abs/2107.01038v15. Full rights notice: licenses/arxiv-2107.01038-CC-BY-4.0.txt.\par\smallskip

\noindent\textit{Editorial scope.} Counted unit: the Lem whose source label is \texttt{lem: common squarefree} in the article (source0.tex lines 580--582, source label \texttt{lem: common squarefree}), delivered together with the article's own complete original proof of it (source0.tex lines 583--634, one \texttt{proof} environment of 52 lines). No mathematics is written, repaired or re-derived: statement and proof are the article's own text line for line. Required context retained uncounted: eq: associativity, d (lines 454--458); eq: associativity, u (lines 454--458); eq: quadrilateral decomposition (lines 460--464); eq: radical expression (lines 488--490); lem: observable sets with 0 are algebraic (lines 519--522); eq: all 3 terms non-vanishing (lines 526--529). Every definition, display and label of those lines is retained unchanged. Omitted from this package and not redistributed: 1--453, 459--459, 465--487, 491--518, 523--525, 530--579, 635--1517. Nothing inside a retained excerpt is omitted or altered: every retained body is delivered exactly as the article's lines are written, so no omission receipt of any kind accompanies this package. Citations: the retained excerpts carry no citation command at all, so no bibliography entry of the article is transcribed and no reference is redistributed. Reference adaptation: none was needed and none was made; every reference inside the delivered unit resolves inside it, so no unresolved reference or citation is produced. Printed slips kept literal: no printed word, symbol, hypothesis or number of the counted statement or of its proof was altered. Layout and class: the article's own class is \texttt{article} with packages including amsmath, amsthm, amssymb, flexisym, breqn, lmodern, avant, courier, fontenc, geometry, mathtools, float, appendix, babel; the wrapper is the lane's pinned \texttt{amsart} editorial wrapper at 10pt on A4, which restates the theorem-like environments the retained text needs and adds the article's own macro definitions that the retained text needs (none), with extra wrapper packages (bm, bbm, dsfont, xspace, cancel, mathtools). The delivered numbering is the unit's own. Assets: no retained span contains a figure environment, a graphics-inclusion command, a PDF-page inclusion, a table or a TikZ picture; no image file is copied into this package, the package declares no asset dependency and no image file is redistributed with it. Licence: version arXiv:2107.01038v15 of this article is distributed under the Creative Commons Attribution 4.0 International licence, as recorded in the arXiv licence metadata for this exact version and shown on the abstract page https://arxiv.org/abs/2107.01038v15. A full rights notice is delivered at licenses/arxiv-2107.01038-CC-BY-4.0.txt. Characters: every retained excerpt is pure ASCII, so the delivered engine, XeLaTeX with its default Latin Modern text font, reports no missing character.
\begin{eqnarray}
Y_{\alpha\beta}^{ij}Y_{\beta\gamma}^{ij} & = & -Y_{\alpha\gamma}^{ij},\label{eq: associativity, d}\\
Y_{\alpha\beta}^{im}\cdot Y_{\alpha\beta}^{mj} & = & -Y_{\alpha\beta}^{ij}.
\label{eq: associativity, u}
\end{eqnarray}

\begin{equation}
    Y_{\alpha\beta}^{ij} = -Y_{\alpha\delta}^{ij}\cdot Y_{\delta\beta}^{ij} 
    = -Y_{\alpha\delta}^{im}\cdot Y_{\alpha\delta}^{mj}\cdot Y_{\delta\beta}^{im}\cdot Y_{\delta\beta}^{mj},\quad i,j,m\in\mathcal{I},\,\alpha,\beta,\delta\in\mathcal{I}^{\mathtt{C}}.
    \label{eq: quadrilateral decomposition}
\end{equation}

 \begin{equation}
Y_{\alpha\beta}^{ij}=\frac{\varepsilon_{\alpha\beta}^{ij}\cdot\sqrt{B_{\alpha\beta}^{ij}}+A_{\alpha\beta}^{ij}}{2\cdot h(\mathcal{I}_{\beta}^{i})\cdot h(\mathcal{I}_{\alpha}^{j})}\label{eq: radical expression}
 \end{equation}

\begin{lem}
\label{lem: observable sets with 0 are algebraic}
Any radical term $Y_{\alpha\beta}^{ij}$ forces $0\notin\chi(\mathcal{I}\mid_{\alpha\gamma}^{ij})$ for all $\gamma\in\mathcal{I}^{\mathtt{C}}$, $\gamma\neq \alpha$. 
\end{lem}

 \begin{equation}
h(\mathcal{I})\cdot h(\mathcal{I}_{\alpha\beta}^{ij})\cdot h(\mathcal{I}_{\alpha}^{i})\cdot h(\mathcal{I}_{\beta}^{i})\cdot h(\mathcal{I}_{\alpha}^{j})\cdot h(\mathcal{I}_{\beta}^{j})\neq0
\label{eq: all 3 terms non-vanishing}
 \end{equation}

\begin{lem}
\label{lem: common squarefree} All observable $Y$-terms with basis $\mathcal{I}$ lie in the same quadratic extension of $\mathbb{F}$. 
\end{lem}

\begin{proof}
We can assume the existence of a radical term $Y_{\alpha\beta}^{ij}$, otherwise all the observable $Y$-terms lie in $\mathbb{F}$ and the thesis follows. 
By Lemma \ref{lem: observable sets with 0 are algebraic}, we can focus on observable sets that satisfy (\ref{eq: all 3 terms non-vanishing}) for the rest of the proof: being $\mathbb{C}(\mathbf{t})$ a unique factorisation domain, we can express 
 \begin{equation}
B_{\alpha\beta}^{ij}:=\left(Q_{\alpha\beta}^{ij}\right)^{2}\cdot D\label{eq: factorization square and squarefree}
 \end{equation}
where $Q_{\alpha\beta}^{ij},D\in\mathbb{C}(\mathbf{t})$ and $D$ is a squarefree polynomial, that is, there exists no $P\in\mathbb{C}(\mathbf{t})$ such that $P^{2}\mid D$. Then, we prove   
 \begin{equation}
Y_{\gamma\delta}^{ml}\in\mathbb{F}(\sqrt{Y_{\alpha\beta}^{ij}})=\mathbb{F}(\sqrt{D})
\label{eq: same algebraic extension}
 \end{equation}
for observable terms $Y_{\gamma\delta}^{ml}$, where $\mathbb{F}(\sqrt{Y_{\alpha\beta}^{ij}})$ denotes the algebraic extension of $\mathbb{F}$ by a square root of $Y_{\alpha\beta}^{ij}$.

We start from the instance $\{m,l\}=\{i,j\}$ and $\delta\in\{\alpha,\beta\}$, noting that we only have to consider radical triples $(Y_{\alpha\beta}^{ij},Y_{\beta\gamma}^{ij},Y_{\alpha\gamma}^{ij})$: being (\ref{eq: all 3 terms non-vanishing}) verified by $\chi(\mathcal{I}|_{\alpha\beta}^{ij})$ and $\chi(\mathcal{I}|_{\gamma\delta}^{ij})$, all the components of this triple are observable, and from (\ref{eq: associativity, d}) we get the thesis (\ref{eq: same algebraic extension}) when $Y_{\gamma\alpha}^{ij}\in\mathbb{F}$ or $Y_{\gamma\beta}^{ij}\in\mathbb{F}$. Therefore, taking into account (\ref{eq: radical expression}) for 
%the radical terms $Y_{\alpha\beta}^{ij}$, $Y_{\alpha\gamma}^{ij}$, and $Y_{\gamma\beta}^{ij}$ 
the $Y$-terms in the expression $Y_{\alpha\beta}^{ij}=-Y_{\alpha\gamma}^{ij}\cdot Y_{\gamma\beta}^{ij}$, we get  
%(\ref{eq: three degree-2 terms, expanded}), 
 \begin{equation}
Y_{\alpha\gamma}^{ij}\in\mathbb{F}(\sqrt{B_{\alpha\beta}^{ij}})\Leftrightarrow Y_{\beta\gamma}^{ij}\in\mathbb{F}(\sqrt{B_{\alpha\beta}^{ij}}).
\label{eq: conditionals same extension}
 \end{equation}
Using again the assumption $Y_{\alpha\beta}^{ij}$, $Y_{\alpha\gamma}^{ij}$, $Y_{\gamma\beta}^{ij}\notin\mathbb{F}$, the product $Y_{\alpha\gamma}^{ij}\cdot Y_{\gamma\beta}^{ij}$ lies in a quadratic extension of $\mathbb{F}$ only if (\ref{eq: same algebraic extension}) holds. 
%at $\delta\in\{\alpha,\beta\}$. 
This argument can be adapted to triples $(Y_{\alpha\beta}^{ij},Y_{\alpha\beta}^{im},Y_{\alpha\beta}^{jm})$ by transposition of the upper and lower indices. Consequently, for all $\gamma_{1},\gamma_{2}\in\mathcal{I}^{\mathtt{C}}$ and $m_{1},m_{2}\in\mathcal{I}$, such that $Y_{\gamma_{1}\gamma_{2}}^{ij}$ and $Y_{\alpha\beta}^{m_{1}m_{2}}$ are observable, $Y_{\gamma_{w}\alpha}^{ij}$ and $Y_{\alpha\beta}^{m_{u}i}$ are also observable for all $u,w\in[2]$, and we have 
\begin{equation}
    Y_{\gamma_{1}\gamma_{2}}^{ij}=-Y_{\gamma_{1}\alpha}^{ij}\cdot Y_{\alpha\gamma_{2}}^{ij}\in \mathbb{F}(\sqrt{D}),\quad  Y_{\alpha\beta}^{m_{1}m_{2}}=-Y_{\alpha\beta}^{m_{1}i}\cdot Y_{\alpha\beta}^{im_{2}}\in \mathbb{F}(\sqrt{D}).
\label{eq: same quadratic extension, only upper or lower}
\end{equation}

Then, we move to radical terms $Y_{\gamma_{1}\gamma_{2}}^{m_{1}m_{2}}$ with $\{\alpha,\beta\}\neq\{\gamma_{1},\gamma_{2}\}$ and $\{i,j\}\neq\{m_{1},m_{2}\}$: for all $a,b,u,w\in[2]$, $i_{a}\in\{i,j\}$, and $\alpha_{b}\in\{\alpha,\beta\}$, the set $\chi(\mathcal{I}|_{\gamma_{w}\alpha_{b}}^{m_{u}i_{a}})$ is observable since $h(\mathcal{I}_{\alpha_{b}}^{i_{a}})\cdot h(\mathcal{I}_{\gamma_{w}}^{m_{u}})\neq0$. For any $u,w\in[2]$, we use such sets and (\ref{eq: quadrilateral decomposition}) to express  
\begin{equation} 
Y_{\alpha\beta}^{ij}=-Y_{\alpha\gamma_{w}}^{im_{u}}Y_{\alpha\gamma_{w}}^{m_{u}j}Y_{\gamma_{w}\beta}^{im_{u}}Y_{\gamma_{w}\beta}^{m_{u}j}\notin\mathbb{F}
\end{equation} 
which implies that there exists at least one pair in $\{i,j\}\times\{\alpha,\beta\}$, say $(i,\alpha)$ with proper labelling, such that $Y_{\alpha\gamma_{w}}^{im_{u}}$ is radical. This condition forces $h(\mathcal{I}_{\gamma_{w}}^{j})\cdot h(\mathcal{I}_{\beta}^{m_{u}})\neq0$ too: indeed, at $h(\mathcal{I}_{\gamma_{w}}^{j})=0$ we would have $Y_{\alpha\gamma_{w}}^{jm_{u}},Y_{\alpha\gamma_{w}}^{ij}\in\mathbb{F}$, since they derive from observable sets containing $0$, so their product would return $Y_{\alpha\gamma_{w}}^{im_{u}}\in\mathbb{F}$ by (\ref{eq: associativity, u}). A similar argument gives $h(\mathcal{I}_{\beta}^{m_{u}})\neq0$. In this way, we get 
\begin{equation}
h(\mathcal{I}_{\gamma_{w}}^{i_{a}})\cdot h(\mathcal{I}_{\alpha_{b}}^{m_{u}})\neq0,
\quad a,b,u,w\in[2],\,i_{a}\in\{i,j\},\,\alpha_{b}\in\{\alpha,\beta\}.
\label{eq: non-central non-vanishing}
\end{equation}

%We show that this configuration generates a set of radical sets connecting $\chi(\mathcal{I}\mid_{\alpha\beta}^{ij})$ and $\chi(\mathcal{I}\mid_{\gamma_{1}\gamma_{2}}^{m_{1}m_{2}})$. 
When $\chi(\mathcal{I}\mid_{\alpha\beta}^{m_{1}m_{2}})$ is radical, we apply twice (\ref{eq: same quadratic extension, only upper or lower}) to obtain 
\begin{equation} 
Y_{\alpha\beta}^{ij}\in\mathbb{F}(\sqrt{D})\setminus \mathbb{F} \Rightarrow Y_{\alpha\beta}^{m_{1}m_{2}}\in\mathbb{F}(\sqrt{D})\setminus\mathbb{F} \Rightarrow Y_{\gamma_{1}\gamma_{2}}^{m_{1}m_{2}}\in\mathbb{F}(\sqrt{D})\setminus\mathbb{F}
\end{equation} 
and a similar implication holds for radical $\chi(\mathcal{I}\mid_{\gamma_{1}\gamma_{2}}^{ij})$. When both $\chi(\mathcal{I}\mid_{\alpha\beta}^{m_{1}m_{2}})$ and $\chi(\mathcal{I}\mid_{\gamma_{1}\gamma_{2}}^{ij})$ contain $0$, (\ref{eq: non-central non-vanishing}) entails $h(\mathcal{I}_{\alpha\beta}^{m_{1}m_{2}})=h(\mathcal{I}_{\gamma_{1}\gamma_{2}}^{ij})=0$; being $Y_{\alpha\beta}^{ij},Y_{\gamma_{1}\gamma_{2}}^{m_{1}m_{2}}\notin\mathbb{F}$, this means that we can find $\overline{i}\in\{i,j\}$ and $\overline{m}\in\{m_{1},m_{2}\}$ such that $Y_{\alpha\beta}^{\overline{i}m_{u}}$ and $Y_{\omega_{1}\omega_{2}}^{\overline{m}x}$ are radical for all $u\in[2]$ and $x\in\{i,j\}$. In particular, both $Y_{\alpha\beta}^{\overline{i}\overline{m}}$ and $Y_{\omega_{1}\omega_{2}}^{\overline{i}\overline{m}}$ are radical: so, we apply (\ref{eq: same quadratic extension, only upper or lower}) three times to get 
\begin{eqnarray}
    Y_{\alpha\beta}^{ij}\in\mathbb{F}(\sqrt{D})\setminus \mathbb{F} & \Rightarrow & 
    Y_{\alpha\beta}^{\overline{i}\overline{m}}\in\mathbb{F}(\sqrt{D})\setminus\mathbb{F} \nonumber \\ 
    & \Rightarrow &  Y_{\omega_{1}\omega_{2}}^{\overline{i}\overline{m}}\in\mathbb{F}(\sqrt{D})\setminus\mathbb{F} \Rightarrow Y_{\omega_{1}\omega_{2}}^{m_{1}m_{2}}\in\mathbb{F}(\sqrt{D})\setminus\mathbb{F}.
\end{eqnarray}
\end{proof}

\end{document}
